%%%%%%%%%%%%%%%%%%%%%%%%%%%%%%%%%%%%%%%%%
% Medium Length Professional CV
% LaTeX Template
% Version 3.0 (December 17, 2022)
%
% This template originates from:
% https://www.LaTeXTemplates.com
%
% Author:
% Vel (vel@latextemplates.com)
%
% Original author:
% Trey Hunner (http://www.treyhunner.com/)
%
% License:
% CC BY-NC-SA 4.0 (https://creativecommons.org/licenses/by-nc-sa/4.0/)
%
%%%%%%%%%%%%%%%%%%%%%%%%%%%%%%%%%%%%%%%%%

%----------------------------------------------------------------------------------------
%	PACKAGES AND OTHER DOCUMENT CONFIGURATIONS
%----------------------------------------------------------------------------------------

\documentclass[
	%a4paper, % Uncomment for A4 paper size (default is US letter)
	11pt, % Default font size, can use 10pt, 11pt or 12pt
]{resume} % Use the resume class

% in your .tex preamble (after \documentclass)
\usepackage{tikz}
\usepackage{enumitem}
\setlength{\emergencystretch}{2em}
\setlength{\medskipamount}{4pt plus 2pt minus 1pt}

\newsavebox{\linkboxreference}

\newcommand{\linkbox}[2]{%
  \sbox{\linkboxreference}{Google Scholar}%
  \tikz[baseline=(X.base)]\node[
    inner xsep=4pt, inner ysep=2pt,
    rounded corners=2pt,
    draw=cyan!70, line width=0.6pt
  ] (X) {\makebox[\wd\linkboxreference][c]{\vphantom{Google Scholar}\href{#1}{#2}}};%
}


%------------------------------------------------

\name{Yunwoo Lee} % Your name to appear at the top

% You can use the \address command up to 3 times for 3 different addresses or pieces of contact information
% Any new lines (\\) you use in the \address commands will be converted to symbols, so each address will appear as a single line.
\address{
    Phone: (206)~$\cdot$~843~$\cdot$~4857 \\
    Email: \href{mailto:yunu919@uw.edu}{yunu919@uw.edu}
}

\address{Seattle, WA 98105}

\address{
  \linkbox{https://yunlee919.github.io/}{Website}
  \linkbox{https://scholar.google.com/citations?user=vN7gwKEAAAAJ}{Google Scholar}
  \linkbox{https://github.com/yunlee919}{GitHub}
}



%----------------------------------------------------------------------------------------

\begin{document}

%----------------------------------------------------------------------------------------
%	EDUCATION SECTION
%----------------------------------------------------------------------------------------

\begin{rSection}{Education}

	\textbf{University of Washington} \hfill Seattle, WA \\
	\textsl{Master of Science in Computational Linguistics} \hfill
    \textsl{Sep. 2025 -- Present}

    \textbf{\href{https://www.hufs.ac.kr/hufs/index.do\#section1}{Hankuk University of Foreign Studies, HUFS}} \hfill Seoul, South Korea \\
    \textsl{Bachelor of Arts in Linguistics \& Cognitive Science} \hfill
    \textsl{Mar. 2018 -- Feb. 2024} \\
	\textsl{Bachelor of Language Science in Artificial Intelligence (Double Major)}\\
    Advisor: Jeesun Nam
\end{rSection}

%----------------------------------------------------------------------------------------
%	RESEARCH EXPERIENCE SECTION
%----------------------------------------------------------------------------------------

\begin{rSection}{Research Experience}
\begin{rSubsection}{\href{https://lt4cpr.github.io/aacl2026-workshop-LT4CPR/index.html}{LT4CPR, UW}}{Seattle, WA}{Research Assistant, Supervisor: Prof. \href{https://faculty.washington.edu/fxia/}{Fei Xia}}{Jun. 2026 -- Present}
    \item Developed evaluation methods for comparing LLM-generated crisis situation reports with reference reports, focusing on semantic content preservation and hierarchical discourse alignment
    \item Designed document-, section-, subsection-, and bullet-level evaluation pipelines combining lexical and semantic similarity metrics with bipartite matching for fine-grained content alignment
\end{rSubsection}
%------------------------------------------------

\begin{rSubsection}{\href{http://dicora.kr/}{DICORA Computational Linguistics Lab, HUFS}}{Seongnam, South Korea}{Undergraduate Research Intern, Supervisor: Prof. Jeesun Nam}{Jan. 2022 -- Sep. 2022}
    \item Conducted corpus-based research on Korean-language data for chatbot NLU in financial-news sentiment analysis, focusing on Korean-specific morphosyntactic patterns involving verbal endings, case particles, and object expressions
    \item Contributed to corpus construction and linguistic annotation based on these patterns for chatbot training
\end{rSubsection}

\end{rSection}

%----------------------------------------------------------------------------------------
%	PUBLICATIONS SECTION
%----------------------------------------------------------------------------------------

\begin{rSection}{Publications}
\begin{enumerate}[label={[\arabic*]}, leftmargin=*, itemsep=0.5em, topsep=0pt]
\item \textit{Overview of the LT4CPR 2026 Shared Task: Generating Situation Reports from Social Media Messages.} Fei Xia, Haotian Zhu, \textbf{Yunwoo Lee}, Eitan Grenev, Allison K. Burton, Yingzhe Fan, William D. Lewis, Antonis Anastasopoulos, and Belu Ticona.\\
\textit{AACL-IJCNLP Workshop on Language Technologies for Crisis Preparedness and Response (LT4CPR),}\\
2026. (Accepted)

\item \textit{Relational Grounding Failures in Zero-Shot Cross-Lingual Dialogue Summarization with Small Language Models.}
\textbf{Yunwoo Lee}, Ziwa Li, and Jennifer Huang.
\textit{EMNLP Workshop on Grounding Language Models: Learning Faithfully and Efficiently (GroundLM),} 2026. (Accepted)
\end{enumerate}

\end{rSection}

%----------------------------------------------------------------------------------------
%	INDUSTRY EXPERIENCE SECTION
%----------------------------------------------------------------------------------------

\begin{rSection}{Industry Experience}

	\begin{rSubsection}{\href{https://www.keti.re.kr/eng/main/main.php}{Korea Electronics Technology Institute (KETI)}}{Seongnam, South Korea}{Researcher, Language Team - AIRC (Artificial Intelligence Research Center)}{Sep. 2024 -- Jun. 2025}
		\item Developed trustworthiness benchmarks for Korean LLMs (hallucination, reliability, sociocultural bias) through prompt design, rubrics, and failure-mode analysis
        \item Validated a Korean multimodal dialogue dataset by aligning utterance-level pragmatic functions with nonverbal cue labels to improve human–AI interaction robustness
	\end{rSubsection}

%------------------------------------------------

	\begin{rSubsection}{\href{https://nc-ai.com/en}{NCSOFT}}{Seongnam, South Korea}{Language AI Researcher, Language Data Team}{Mar. 2024 -- Sep. 2024}
		\item Led AI Red Team initiatives, creating robust testing protocols and diverse conversational datasets to identify and mitigate ethical and safety risks in language models
		\item Categorized language model vulnerabilities into single-turn versus multi-turn threats by identifying discourse-level risks, including anaphoric references that are undetectable within a single turn
	\end{rSubsection}

%------------------------------------------------

	\begin{rSubsection}{\href{https://www.sktelecom.com/index_en.html}{SK TELECOM}}{Seoul, South Korea}{Linguistic Data Analyst, AI Technology Unit}{Feb. 2023 -- Jun. 2023}
		\item Led human-centric annotation to ensure sociolinguistic coverage (honorific levels, informal slang, code-mixing) in conversational training data
		\item Created a linguistically grounded ambiguity taxonomy and applied it to normalize Korean queries and improve chatbot robustness
	\end{rSubsection}

\end{rSection}


%----------------------------------------------------------------------------------------
%	AWARDS AND HONORS SECTION
%----------------------------------------------------------------------------------------

\begin{rSection}{Awards and Honors}
    \textbf{Graduate Student Conference Presentation Award, UW} \hfill 2026 \\
    \textbf{\href{https://linguistics.washington.edu/honors-awards\#S06}{CLMS Scholarship, UW}} (\$14,500) \hfill 2025 \\
    \textbf{Departmental Scholarship, HUFS} \hfill 2021, 2023 \\
    \textbf{Outstanding Undergraduate Thesis Award, HUFS} \hfill 2022 \\
    \textbf{Best Composition Award, HUFS} \hfill 2018 \\
    \textbf{Student Leadership Scholarship, HUFS} \hfill 2018 \\
    \textbf{National Merit Scholarship} \hfill 2018
\end{rSection}

%----------------------------------------------------------------------------------------
%	PROFESSIONAL SERVICE SECTION
%----------------------------------------------------------------------------------------

\begin{rSection}{Professional Service}
\begin{enumerate}[label={[\arabic*]}, leftmargin=*, itemsep=0.5em, topsep=0pt]
\item Reviewer, \textbf{\href{https://vlm4rwd.github.io/}{VLM4RWD 2026 Workshop}}, NeurIPS 2026
\end{enumerate}

\end{rSection}

%----------------------------------------------------------------------------------------
%	RESEARCH PROJECTS SECTION
%----------------------------------------------------------------------------------------



\begin{rSection}{Research Projects}
    \begin{rSubsection}{Learning Linguistic Generalizations Across Languages from Indirect Evidence: A Filtered-Corpus Study of English and Korean}{}{M.S. Thesis, University of Washington}{Oct. 2026 -- Present}
        \item Investigating cross-linguistic grammatical generalization in language models through bidirectional English--Korean filtered-corpus training experiments.
        \item Evaluating whether construction-specific evidence from one language facilitates generalization to unseen grammatical constructions in another language using minimal-pair grammaticality judgments.
    \end{rSubsection}


	\begin{rSubsection}{Linguistic Pattern Analysis for Financial Sentiment: Attribute-Verb Relationships in Stock Market Articles}{}{B.A. Thesis, HUFS}{Sep. 2022 -- Dec. 2022}
		\item Modeled semantic constraints between attribute nouns (e.g., interest rates, costs) and directional predicates (rise/fall) to capture context-dependent polarity shifts in financial discourse
		\item Recognized with the Outstanding Undergraduate Thesis Award
	\end{rSubsection}
\end{rSection}

%----------------------------------------------------------------------------------------
%	LANGUAGES AND TECHNICAL SKILLS SECTION
%----------------------------------------------------------------------------------------

\begin{rSection}{Languages and Technical Skills}

	\begin{tabular}{@{} >{\bfseries}l @{\hspace{6ex}} l @{}}
		Programming & Python, R, C++, SQL, Java \\
        ML/NLP & PyTorch, NLTK \\
        Tools & Git, Linux/CLI \\
        Typesetting & \LaTeX \\
		Languages & Korean (native), English (fluent) \\
	\end{tabular}

\end{rSection}

%----------------------------------------------------------------------------------------
%	EXAMPLE SECTION
%----------------------------------------------------------------------------------------

%\begin{rSection}{Section Name}

	%Section content\ldots

%\end{rSection}

%----------------------------------------------------------------------------------------

\end{document}
